Cosmological redshift underpins myriad applications in modern astronomy. Accurate redshift estimates permit the conversion of observed magnitudes of distant galaxies into intrinsic luminosities, allowing studies of the galaxy luminosity function \citep[e.g.][]{1976ApJ...203..297S, 2015ApJ...810...71F, 2023MNRAS.523.1009B}, the evolution of galaxy size and morphology across cosmic epochs \citep[e.g.][]{2014ARA&A..52..291C, 2015ApJS..219...15S}, and scaling relations between galaxy stellar masses and properties of their interstellar medium \citep[e.g.][]{1998ApJ...498..541K, 2004ApJ...613..898T}, central supermassive black holes \citep[e.g.][]{2013ARA&A..51..511K}, and host haloes \citep[e.g.][]{2013ApJ...770...57B, 2018ARA&A..56..435W}.

In cosmology, Type Ia supernova host galaxy redshifts provided the first direct evidence of the accelerating expansion of the Universe \citep{1998AJ....116.1009R, 1999ApJ...517..565P}. Baryon Acoustic Oscillations (BAOs), measured from galaxy clustering, act as a standard ruler for cosmic distances \citep{2005ApJ...633..560E}, while large-scale structure surveys map the growth of density fluctuations \citep{1980lssu.book.....P, 2006Natur.440.1137S}. The weak gravitational lensing of distant galaxies \citep{1990ApJ...349L...1T, 1998ApJ...498...26K, 2001PhR...340..291B} enables the reconstruction of the matter distribution and serves as a sensitive probe of dark energy through cosmic-shear tomography \citep{2015RPPh...78h6901K}. Collectively, these observables demonstrate how cosmological redshift is the keystone linking local galaxy properties to the geometry and dynamics of the Universe, and a critical underpinning of the concordance \(\Lambda\)CDM model.

Recent wide-field imaging surveys offer multi-wavelength photometry for hundreds of millions of galaxies, such as the Kilo-Degree Survey \citep[KiDS\footnote{\url{https://kids.strw.leidenuniv.nl}};][]{2024A&A...686A.170W}, the Dark Energy Survey \citep[DES\footnote{\url{https://www.darkenergysurvey.org}};][]{2025arXiv250105739B}, and the Hyper Suprime-Cam Survey \citep[HSC\footnote{\url{https://hsc.mtk.nao.ac.jp/ssp/}};][]{2022PASJ...74..421L}. Nonetheless, acquiring precise redshifts using spectroscopy for such extensive samples remains exceedingly resource-demanding. As a result, studies on large-scale structures and galaxy evolution are increasingly dependent on photometric redshift methods, which estimate redshifts from broadband flux observations, a technique initially proposed by \citet{1985AJ.....90..418K, 1986ApJ...303..154L}.

Next‐generation facilities -- including the NSF-DOE Vera C. Rubin Observatory \citep[\textit{Rubin}\footnote{\url{https://rubinobservatory.org}};][]{2009arXiv0912.0201L}, the Euclid mission \citep[\textit{Euclid}\footnote{\url{https://www.cosmos.esa.int/web/euclid}};][]{2025A&A...697A...1E}, the Nancy Grace Roman Space Telescope \citep[\textit{Roman}\footnote{\url{https://roman.gsfc.nasa.gov}};][]{2015arXiv150303757S} and the China Space Station Telescope \citep[CSST\footnote{\url{http://www.bao.ac.cn/csst/}};][]{2019ApJ...883..203G} -- will deliver unprecedented imaging depth and sky coverage. Exploiting their full cosmological potential to probe dark matter and dark energy, however, demands exceptional control of systematic uncertainties, especially in photometric redshift performance \citep{2018ARA&A..56..393M}.

In recent decades, numerous techniques for estimating photometric redshifts have been developed and are generally classified into two categories based on the type of data and the priors used during inference, as reviewed in \citet{2022ARA&A..60..363N}. Comprehensive descriptions of most of these techniques are available through the Redshift Assessment Infrastructure Layers \citep[\texttt{RAIL}\footnote{\url{https://github.com/LSSTDESC/rail}};][]{2025arXiv250502928T} platform for Legacy Survey of Space and Time (LSST) Dark Energy Science Collaboration (DESC). Specifically, the most traditional approach, template fitting, matches observed broadband photometry to libraries of galaxy Spectral Energy Distribution (SED) templates, using software packages such as \texttt{LePhare} \citep{1999MNRAS.310..540A,2006A&A...457..841I}, \texttt{BPZ} \citep{2000ApJ...536..571B, 2006AJ....132..926C}, \texttt{ZEBRA} \citep{2006MNRAS.372..565F}, \texttt{EAZY} \citep{2008ApJ...686.1503B}, and \texttt{DELIGHT} \citep{2017ApJ...838....5L}. 

Modern template-based approaches incorporate sophisticated Stellar Population Synthesis (SPS) models -- such as \texttt{FSPS} \citep{2009ApJ...699..486C, 2010ApJ...708...58C, 2010ApJ...712..833C}, \texttt{BAGPIPES} \citep{2018MNRAS.480.4379C}, and \texttt{Prospector} \citep{2021ApJS..254...22J} -- that include flexible star-formation histories and improved dust and nebular-emission treatments. These developments have enabled highly competitive photometric-redshift point estimates and ensemble redshift distributions, particularly for surveys with medium- or narrow-band photometry where spectral features are more cleanly resolved \citep[e.g.][]{2009ApJ...690.1250S, 2011ApJ...742...61S, 2021MNRAS.501.6103A}. For stage-IV, broad-band imaging surveys, however, the dominant systematic uncertainties generally arise from the realism and flexibility of the adopted galaxy-population priors rather than from intrinsic limitations of the SPS models themselves. Although template fitting provides a principled route to inferring redshift posteriors from noisy broad-band data, mismatches between the assumed and true population priors can still introduce biases that ultimately limit the accuracy of broad-band photometric-redshift estimates \citep[e.g.][]{2010A&A...523A..31H,2013ApJ...775...93D}.

Conversely, a variety of machine-learning algorithms have been devised to characterise the complex relationship between galaxy broadband SEDs and redshifts, utilising galaxies with known redshifts. These include artificial neural networks \citep[\texttt{ANNz};][]{2004PASP..116..345C, 2016PASP..128j4502S}, decision trees with random forest techniques \citep[\texttt{TPZ};][]{2013MNRAS.432.1483C}, Gaussian processes \citep[\texttt{GPz};][]{2016MNRAS.462..726A}, the quasi-Newton algorithm \citep[\texttt{MLPQNA};][]{2012A&A...546A..13C}, hybrid density estimation \citep[\texttt{METAPHOR};][]{2017MNRAS.465.1959C}, boosted decision trees combined with basis function decomposition \citep[\texttt{FlexZBoost};][]{10.1214/17-EJS1302, 2020A&C....3000362D, 2020MNRAS.499.1587S}, the normalising flow algorithm \citep[\texttt{PZFlow}; ][]{2024AJ....168...80C}, and stratified learning \citep[\texttt{StratLearn-z};][]{2018MNRAS.473.3969R, 2024MNRAS.534.3808A, 2025OJAp....8E..50M}. 

These data-driven approaches rely on training samples with high completeness and precisely measured redshifts, typically drawn from wide‐area spectroscopic surveys such as the Sloan Digital Sky Survey \cite[SDSS\footnote{\url{https://www.sdss.org}};][]{2000AJ....120.1579Y}, the Two-degree Field Galaxy Redshift Survey \citep[2DFGRS\footnote{\url{http://www.2dfgrs.net}};][]{2001MNRAS.328.1039C} and the Galaxy And Mass Assembly survey \citep[GAMA\footnote{\url{https://www.gama-survey.org}};][]{2011MNRAS.413..971D}, as well as deep spectroscopic programmes (e.g.\ zCOSMOS; \citealp{2007ApJS..172...70L}, DEEP2; \citealp{2013ApJS..208....5N}, VIPERS; \citealp{2014A&A...566A.108G}). Narrow‐band imaging campaigns (e.g.\ J‐PAS\footnote{\url{https://www.j-pas.org}}; \citealp{2014arXiv1403.5237B}, PAUS\footnote{\url{https://www.paus-survey.org}}; \citealp{2025A&A...693A.102D}) extend this reach to fainter magnitudes and higher redshifts. New and forthcoming facilities -- among them the Dark Energy Spectroscopic Instrument \citep[DESI\footnote{\url{https://www.desi.lbl.gov}};][]{2024AJ....167...62D}, the Subaru Prime Focus Spectrograph \citep[PFS\footnote{\url{https://pfs.ipmu.jp}};][]{2014PASJ...66R...1T}, the 4-metre Multi-Object Spectroscopic Telescope \citep[4MOST\footnote{\url{https://www.4most.eu/cms/home/}};][]{2016SPIE.9910E..1NW} and William Herschel Telescope Enhanced Area Velocity Explorer \citep[WEAVE\footnote{\url{https://www.ing.iac.es/astronomy/instruments/weave/weaveinst.html}};][]{2024MNRAS.530.2688J} -- promise substantially improved spectroscopic completeness, thereby enhancing the training and accuracy of machine‐learning photometric‐redshift estimators.

Nevertheless, spectroscopic samples remain subject to magnitude limits and exhibit uneven coverage in redshift-colour space, resulting in non‐representative training datasets and inducing selection biases \citep{2020A&A...637A.100W, 2020MNRAS.496.4769H}. Spectroscopic misidentifications and catalogue mismatches further introduce systematic offsets in point estimates and distort the shapes of posterior redshift distributions \citep{2015ApJ...813...53M, 2020MNRAS.496.4769H}. Moreover, degeneracies in galaxy SEDs and redshift-colour ambiguities exacerbate these biases. Mitigating these challenges therefore demands robust statistical frameworks, truly representative calibration samples and explicit correction schemes that account for spectroscopic incompleteness.

To address these limitations, an alternative method introduced by \citep{2024ApJ...967L...6M} utilises synthetic mock galaxy catalogues to supplement missing galaxy SEDs, thereby providing more representative datasets for training machine learning methods to estimate photometric redshifts. Our work is therefore complementary to SPS-based methods: instead of modifying SED priors, we aim to improve the empirical completeness of the training sample in colour-magnitude space through targeted data augmentation. This requires a careful understanding of the realism of the supplemental catalogue, but was shown to provide significant improvements, particularly in high-redshift estimation where the training dataset is least representative. 

In this study, we extend this framework by using a Self-Organising Map (SOM) to identify regions of colour-magnitude space with insufficient spectroscopic-redshift coverage and preferentially augment those regions. This enables a more targeted and effective augmentation of simulated galaxies, leading to improved accuracy in photometric-redshift estimates. We also introduce two-stage controlled variations to generate a broad suite of synthetic spectroscopic samples representative of those potentially available for LSST. This allows us to quantify the impact of differing spectroscopic selection functions and to assess photometric-redshift performance under both one-year (Y1) and ten-year (Y10) LSST depths.

This work focuses on generating mock photometric and spectroscopic galaxy datasets with a realistic level of fidelity for LSST cosmological analyses, and demonstrates how SOM-based data augmentation improves photometric-redshift estimates for individual galaxies. In a forthcoming study (Zhang et al., in preparation), we extend this framework to establish tomographic binning schemes and to calibrate the ensemble redshift distributions for each tomographic bin.

The structure of the paper is as follows: Section~\ref{sec:2} provides a review of the mock galaxy catalogues; Section~\ref{sec:3} outlines our photometric and spectroscopic sample selection and augmentation process; Section~\ref{sec:4} explains our photometric redshift modelling; Section~\ref{sec:5} evaluates the precision achieved; and Section~\ref{sec:6} concludes with a summary of our findings and and presents opportunities for future research.